\section{Integration with the BX3 Framework}
\label{sec:rs-integration}

The Recursive Spawning Protocol is a structural instantiation of the BX3 three-layer accountability architecture, not a module composed with it as a separate component. Every RSP mechanism maps onto exactly one BX3 layer, and this mapping is minimal and structurally necessary: no RSP guarantee can be achieved with fewer than all three layers, and no layer can subsume another's role without collapsing at least one correctness property. This section specifies the three mappings and demonstrates why each is exclusive to its layer.

% -------------------------------------------------------
\subsection{Purpose Layer: Accountability Anchor}
\label{sec:rs-purpose-anchor}

The Purpose Layer anchors RSP along two structurally singular axes that no other layer can occupy.

\medskip
\noindent\textbf{Exclusive origin of spawn authorization.} At root provisioning, the Purpose Layer defines and records the spawn authorization policy $P_{\mathrm{spawn}}$ in the Fact Layer authorization ledger. $P_{\mathrm{spawn}}$ encodes two mandatory pre-conditions for any child spawn: (i) the child operating scope must be a strict descendant refinement of the parent scope ($R(n_{i+1}) \subset R(n_i)$), and (ii) the spawn must be operationally necessary --- the unmet condition must be unsatisfiable within the parent's existing scope. These conditions originate exclusively at the Purpose Layer. The Bounds Engine evaluates proposed spawns against $P_{\mathrm{spawn}}$; the Fact Layer enforces it; neither originates it. No machine actor may issue a spawn warrant in the absence of a matching Purpose Layer authorization record. The separation is architectural, not organizational: it holds regardless of privilege level, operational urgency, or system state.

\medskip
\noindent\textbf{Exclusive terminus of escalation.} When the chain reaches $D_{\mathrm{ESCALATE}}$, the protocol compulsory-escalates to the human anchor $h(n)$. Critically, $h(n)$ is non-collapsing: for all nodes $n_i$ in chain $C$, $h(n_i) = h(n_0)$. The human anchor does not migrate down the chain as depth increases. The human issues one of three directives --- \texttt{ACK}, \texttt{RESOLVE}, or \texttt{REJECT} --- and no machine actor may issue any of these. The chain terminates in human judgment, never in autonomous resolution.

This is not a policy convention. It is the structural expression of the upstream BX3 accountability guarantee: systems built on the BX3 Framework fail upward into human consciousness. If any machine actor could absorb an exception or issue a terminal directive, the chain would terminate at that actor rather than at a human. The Purpose Layer's exclusivity here is load-bearing: removing it causes the upstream accountability guarantee to collapse into a machine-actor loop.

\medskip
\noindent\textbf{Structural necessity of the Purpose Layer.} The Purpose Layer's two roles are jointly exclusive to it because both require the exercise of intent and the bearing of final accountability --- capabilities structurally reserved for the human principal. The Bounds Engine cannot originate $P_{\mathrm{spawn}}$ because it operates within, not above, the operating range defined by the Purpose Layer; it has no mandate to define what counts as authorized purpose. The Fact Layer cannot serve as the escalation terminus because it is architecturally a deterministic enforcement layer, not a judgment layer: it applies rules, it does not decide. Only the Purpose Layer, operating at the human's behest, can originate intent and bear final accountability. The mapping is therefore structurally necessary, not merely convenient.

% -------------------------------------------------------
\subsection{Bounds Engine: Depth Enforcement}
\label{sec:rs-bounds-depth}

The Bounds Engine provides constrained execution evaluation: it evaluates proposed actions against the operating range defined by the Purpose Layer and enforces the depth constraint that prevents unbounded chain growth. Critically, the Bounds Engine cannot exceed these constraints under any operational circumstances --- there is no privilege level, operational exigency, or system state under which it may override the depth bound.

\medskip
\noindent\textbf{Depth evaluation at spawn.} At each spawn event, the Bounds Engine evaluates current chain depth $|C|$ against the Purpose-Layer-defined maximum $D_{\mathrm{max}}$, both anchored in the Fact Layer authorization ledger. If $|C| \geq D_{\mathrm{max}}$, the Bounds Engine does not execute the spawn. It transitions to \texttt{ESCALATING} state and initiates the Bailout Protocol, propagating the exception upward to $h(n)$. This enforcement is deterministic: the Bounds Engine has no discretion to exceed $D_{\mathrm{max}}$ regardless of urgency, operational exigency, or the importance of the unmet condition. The depth constraint is a hard architectural boundary, not a heuristic or a default that can be overridden.

\medskip
\noindent\textbf{Scope refinement evaluation.} The Bounds Engine also evaluates the scope subset condition $R(n_{i+1}) \subset R(n_i)$ at each spawn event, determining whether a spawn request is justified by operational necessity within the Purpose Layer's authorization scope. This evaluation is a precondition for pre-commit hook invocation; it is not sufficient alone, because the scope refinement condition is independently enforced by the Fact Layer pre-commit hook. The separation reflects the depth-limit insufficiency theorem (Section~\ref{sec:drift-depth-limits}): depth limits alone cannot prevent scope creep drift, so both depth enforcement and scope refinement enforcement are required as structurally distinct operations.

\medskip
\noindent\textbf{Structural reinforcement by the Fact Layer.} The Bounds Engine's depth evaluation is structurally reinforced by the Fact Layer pre-commit hook. Before any spawn operation is recorded, the Fact Layer independently verifies two conditions: (a) $R(n_{i+1}) \subseteq R(n_i)$ (scope subset), and (b) $\mathrm{depth}(n_{i+1}) \leq D_{\mathrm{max}}$ (depth bound). The Fact Layer rejects any spawn violating either condition. The Bounds Engine cannot execute a rejected spawn; it may only transition to \texttt{ESCALATING}. The pre-commit hook makes the depth bound and scope refinement constraint structural invariants --- not policy conventions overridable in exigent circumstances.

The structural necessity of the Bounds Engine follows from the insufficiency of depth limits alone: the Fact Layer can reject a spawn that exceeds the depth bound, but without the Bounds Engine there is no layer responsible for evaluating whether a spawn is operationally necessary within the authorized scope, nor any layer to manage the transition to \texttt{ESCALATING} state when the depth bound is reached. The Fact Layer enforces what the Bounds Engine evaluates; neither layer alone is sufficient.

% -------------------------------------------------------
\subsection{Fact Layer: Forensic Ledger}
\label{sec:rs-fact-ledger}

The Fact Layer provides deterministic enforcement and forensic auditability. It maintains the append-only ledger that makes the RSP chain auditable at any future time and enforces the write protocol constraints that make RSP structurally tamper-evident. The Fact Layer is the structural foundation on which both the Purpose Layer's authorization policy and the Bounds Engine's depth constraints become runtime-enforceable, not merely documented.

\medskip
\noindent\textbf{Append-only forensic ledger.} The Fact Layer records every spawn event, every state transition, every Purpose Layer directive received, and every unresolved exception condition. Each entry is timestamped, attributed to the node that generated it, and hash-chained to the preceding entry. Hash-chaining ensures that no entry can be inserted, deleted, reordered, or altered after recording: any tampering breaks hash continuity and is detectable by any observer with ledger read access. The ledger is the single source of truth for all authorized chain state --- the runtime artifact that distinguishes RSP chains from conventional delegation chains, which exist only as forensic reconstructions attempted after the fact.

The forensic ledger also closes the accountability loop: it records every directive issued by the human anchor and every state transition resulting from that directive, creating a complete, tamper-evident provenance chain from the root human principal to every active node in the system. This is what makes RSP's correctness properties empirically verifiable, not merely analytically sound.

\medskip
\noindent\textbf{Immutable parent pointer.} Once $\alpha(n)$ is established at node provisioning, no actor --- including the Purpose Layer after provisioning --- may modify it. The Fact Layer write protocol rejects any attempt to overwrite $\alpha(n)$. Chain topology changes occur only through authorized spawn operations recorded in the ledger; there is no mechanism for retroactive pointer manipulation. The parent pointer is not merely access-controlled --- it is structurally immutable by protocol definition. This eliminates the re-parenting drift pathway (Failure Mode 3) by making retroactive chain rewriting structurally impossible regardless of the requestor's identity or privilege level.

\medskip
\noindent\textbf{Pre-commit hook.} The pre-commit hook is the enforcement mechanism that makes all other RSP constraints structural. At every spawn event, it verifies conditions (a) and (b) before the spawn is recorded. Rejected spawns are logged; the Bounds Engine does not execute them; the protocol transitions to \texttt{ESCALATING} if the unmet condition persists. The pre-commit hook operates identically on all machine actors; there is no privileged actor capable of bypassing it. The structural immunity from bypass is what distinguishes the pre-commit hook from access-control-based enforcement: even administrative compromise of a machine actor cannot produce a bypass of the pre-commit hook, because the hook is enforced at the Fact Layer write protocol level, not at the actor's discretion.

\medskip
\noindent\textbf{Structural necessity of the Fact Layer.} The Fact Layer's role is structurally necessary because immutability of the parent pointer and the binding force of the pre-commit hook cannot be achieved without a layer that enforces write protocol constraints at the architectural level. Without the Fact Layer write protocol, immutability is a promise rather than a constraint. A mutable parent pointer enables all drift failure modes (Section~\ref{sec:drift-failure-modes}). Without the pre-commit hook, the scope refinement constraint and depth bound are policy documentation --- enforceable until an actor with sufficient privilege decides they are inconvenient. The Fact Layer is the structural enforcement layer that makes the Purpose Layer's policy and the Bounds Engine's constraints operative, not merely documented.

% -------------------------------------------------------
\subsection{Structural Minimality}
\label{sec:rs-minimality}

The three-layer mapping is minimal and sufficient. Removing any layer causes the corresponding guarantee to fail, as shown in Table~\ref{tab:rs-layer-minimality}.

\begin{table}[h]
\begin{center}
\begin{tabular}{p{3.4cm} p{4.6cm}}
\toprule
\textbf{Removed layer} & \textbf{Guarantee that fails} \\
\midrule
Purpose Layer & Exclusive human terminus eliminated; chain can terminate in machine actors; spawn policy becomes self-authorized by machine actors \\
Bounds Engine & Constrained execution evaluation eliminated; no layer to evaluate spawn necessity or manage escalation transition between pre-commit rejection and mandatory escalation \\
Fact Layer & Immutable parent pointer eliminated; retroactive chain manipulation possible; scope refinement and depth bound become policy conventions, not structural invariants \\
\bottomrule
\end{tabular}
\caption{Layer removal consequences for RSP correctness properties.}
\label{tab:rs-layer-minimality}
\end{center}
\end{table}

The three-layer architecture is the minimal structure that makes RSP's correctness properties unconditionally guaranteed. No subset of layers achieves the same guarantees. This is the integration specification for the Recursive Spawning Protocol within the BX3 Framework.
